A tömegközlekedési rendszerek sajnos nem túl népszerűek az utazók körében, egyrészt az autók kényelme miatt, másrészt a tömegközlekedésben jelenlévő problémák miatt. Ezen problémák közé tartoznak a buszok pontatlansága, a zsúfoltság és a hiányos információk. Sok helyen a buszok nem is rendelkeznek menetrenddel, vagy ha igen, azok nem pontosak. Digitalizáció szempontjából is rengeteg hiányosság van a tömegközlekedésben, így ez az egyik legelavultabb szolgáltatások egyike.

Jelenleg már ott tartunk, hogy az emberek többsége okostelefonnal rendelkezik, különböző alkalmazásokat használ és böngészik az interneten. Ezt rengeteg vállalkozás ki is használja, hogy felgyorsítsa és egyszerűsítse a szolgáltatásait. A tömegközlekedésben is egyre több vállalkozás próbálja megoldani a problémákat, de még mindig sok a hiányosság. Ezeket belátva nem is csoda, hogy az emberek inkább az autózást választják, mint a tömegközlekedést.

Egy jól megtervezett és kivitelezett tömegközlekedési rendszer sokat segítene úgy az embereknek, mint a környezetnek. Egyre több problémát okoz ugyanis a globális felmelegedés, amelynek egyik okozója az autók károsanyag kibocsátása. A tömegközlekedés egyik legnagyobb előnye, hogy sok ember szállításával kevesebb szennyező anyag kerül a levegőbe, mint az autózás esetén. Ezen kívül a tömegközlekedés sokkal kisebb területet foglal el, mint az autózás, így a városokban is sokkal kevesebb helyet foglal el a parkolás. Egyes városokban már most is problémát okoz a parkolóhelyek hiánya és a hatalmas dugók, amelyeket az autók okoznak. Ezt mindenki tapasztalja nap mint nap, de sajnos így is kevesen választják a tömegközlekedést.

Azért választottuk ezt a témát Zsigmond Attila kollégámmal, mert mindketten megtapasztaltuk a tömegközlekedés és a városi forgalom problémáit a mindennapi ingázás során és úgy gondoltuk, hogy egy digitalizált rendszer sokat segíthetne a hozzánk hasonló ingázóknak és az alkalmi utazóknak is. Ha az emberek interneten tudnának tájékozódni a buszokról és a forgalomról, intézhetnék itt a jegyvásárlást, és láthatnák a buszok valós idejű pozícióját, valószínűleg sokkal többen választanák a tömegközlekedést az autózás helyett.

A dolgozat célja egy olyan rendszer bemutatása, amely a digitalizáció által megoldja a tömegközlekedésben jelen lévő problémákat, ezáltal kényelmesebbé és egyszerűbbé téve a tömegközlekedést. Ez a rendszer nem csak az utasoknak nyújthat segítséget, hanem a busztársaságoknak is, mert ezt használva egyszerűbben tudják intézni online az ügyeiket és jobban át tudják tekinteni a vállalatukat.

A szóban forgó rendszer neve Bus4U, amely az angol Bus For You rövidítése és amely jelentése magyarul: Buszjárat Neked. A Bus4U egy végfelhasználóknak szánt platformból és egy adminisztrátori felületből tevődik össze. A felhasználói felület főbb funkciói közé tartozik a buszjáratok keresése, a járatokra szóló előzetes online jegyvásárlás, illetve a járatok indulási időpontjainak megtekintése városokra és megállókra lebontva. Ezen kívül megtekinthetőek a buszmegállók a térképen, városok és útvonalak szerint szűrve, valamint valós időben követhetőek az autóbuszok pillanatnyi helyzetei is.

Az adminisztrációs felületen lehetőség nyílik elsősorban a menetrendek és a megállók kezelésére, valamint az útvonalak feltöltésére. A busztársaság adminisztrátora itt tudja kezelni az alkalmazottak információit, és az autóbuszok különböző adatait is, mint az időszakos műszaki vizsga és a biztosítás. Ugyanitt található egy statisztikai felület is, ahol a felhasználók számát, az eladott jegyek számát és az ezekből generált bevételt lehet megtekinteni havi, éves és mindenkori bontásban.

A rendszer két nagy részre osztható fel. Az egyik a backend, amely a háttérfolyamatokért és az üzleti logikáért felelős, míg a másik a frontend, amely a felhasználók számára látható felületet kezeli. A backend működése a kollégám diplomamunkájában kerül ismertetésre \cite{bus4userver}, míg az én dolgozatom a frontend részt mutatja be, amely magában foglalja a weboldalakat és a mobilalkalmazást.

Ebben a dolgozatban szó lesz a frontend követelményeiről, az architektúráról, a felhasználói felület tervezéséről és a megvalósításról. Be fogom mutatni a felhasznált eszközöket és technológiákat, a tervezési és implementációs folyamatot, valamint a végeredményt is. Beszélni fogok a tesztelésről és a továbbfejlesztési lehetőségekről, végül pedig a levonható következtetésekről is.